\subsection{Processo di miglioramento} 
\label{subsec:miglioramento}

\subsubsection{Scopo}
Questo documento definisce le procedure operative per l'analisi, la valutazione e il miglioramento continuo del processo di sviluppo software, al fine di garantire elevati standard di qualità e l'efficienza dei processi.
Lo scopo è quello di definire, controllare e migliorare i processi di \glossario{ciclo di vita} del software.

\subsubsection{Attività}
Le attività operative da eseguire sono:
\begin{enumerate}
    \item definizione dei processi;
    \item valutazione dei processi;
    \item miglioramento dei processi.
\end{enumerate}

\subsubsection{Definizione dei processi}
Vengono analizzati i processi di \glossario{ciclo di vita} del software definiti dallo standard \glossario{ISO\textbackslash IEC 12207:1996} e vengono individuati
i processi necessari alla realizzazione del progetto.
Successivamente viene stabilito il \textit{Way of Warking}, vengono quindi documentate le norme necessarie per implementare correttamente
i processi individuati all'interno del documento \textit{Norme di Progetto}.

\subsubsection{Valutazione dei processi}
Il responsabile di progetto ha il compito di valutare i processi implementati durante l'ultimo \glossario{sprint} di lavoro.
L'obiettivo è quello di individuare eventuali criticità sull'implementazione dei processi e assicurarne la loro efficacia individuando le aree migliorabili.
Per fare questo il responsabile deve inoltre calcolare le \glossario{metriche} definite all'interno del \textit{Piano di Qualifica} che permettono, per determinati processi,
di registrarne l'andamento per effettuare una valutazione più accurata.

\subsubsection{Miglioramento dei processi}
Una volta identificati e valutati i processi implementati durante lo sviluppo e individuati gli aspetti migliorabili nell'implementazione dei processi, è necessario aggiornare le \textit{Norme di Progetto} per avvicinarsi incrementalmente allo stato dell'arte.


\subsubsection{Strumenti}
Gli strumenti elencati in seguito sono approfonditi nella sezione \hyperref[sec:strumenti]{Strumenti}:
\begin{itemize}
    \item \glossario{LaTeX}
    \item \glossario{GitHub}
    \item Discord
\end{itemize}